\documentclass[11pt]{article}
\usepackage[T1]{fontenc}
\usepackage{lmodern,amsmath,amssymb,amsthm}
\usepackage[margin=1in]{geometry}
\usepackage[hidelinks]{hyperref}
\usepackage{microtype}
\hypersetup{pdftitle={BV generators for polynomial Weyl subalgebras},
pdfauthor={Alec Kriebel},
pdfsubject={Ordinary Hochschild cohomology; characteristic-zero squarefree classification},
pdfkeywords={Hochschild cohomology, BV generator, Weyl subalgebra, 30005740}}
\newtheorem{theorem}{Theorem}
\newtheorem{lemma}[theorem]{Lemma}
\newtheorem{proposition}[theorem]{Proposition}
\theoremstyle{remark}\newtheorem{remark}[theorem]{Remark}
\newcommand{\HH}{\operatorname{HH}}
\newcommand{\ad}{\operatorname{ad}}
\newcommand{\Der}{\operatorname{Der}}
\newcommand{\kk}{k}
\renewcommand{\labelitemi}{\(\bullet\)}
\title{BV generators for polynomial Weyl subalgebras}
\author{Alec Kriebel\\
\small Independent researcher\\
\small ORCID: \href{https://orcid.org/0009-0001-9320-500X}{0009-0001-9320-500X}}
\date{22 October 2026\\\small Preprint, version 1.0.0}
\begin{document}
\maketitle
\begin{abstract}
Let \(k\) be any field of characteristic zero and let \(0\ne h\in k[x]\).
For \(A_h=k\langle x,Y\rangle/(Yx-xY-h(x))\), the ordinary Hochschild
cohomology, with its canonical cup product and Gerstenhaber bracket,
admits a Batalin--Vilkovisky generator if and only if \(h\) is squarefree.
In the squarefree case its bracket is zero. A repeated irreducible factor
produces a cup-zero, noncommuting pair of degree-one classes, and a central
annihilator forces every putative generator to vanish on degree one.
Thus in that case no degree-minus-one bracket generator exists, even
without requiring its square to be zero.
The argument uses the published resolution and derivation structure of this
family, and answers the concrete-family question recorded by Lopes in
Oberwolfach Report 46/2023, p.~2702, Question~(2).
\end{abstract}

\section{Question, conventions, and statement}
Throughout, cohomology means ordinary \(\HH^\bullet(A,A)\); its product and
bracket are the canonical ones. A bracket generator is a \(k\)-linear map
\(\Delta:\HH^n\to\HH^{n-1}\) such that, for homogeneous \(a,b\),
\begin{equation}\label{bv}
 [a,b]=(-1)^{|a|}\Delta(a\smile b)
       -(-1)^{|a|}\Delta(a)\smile b-a\smile\Delta(b).
\end{equation}
A BV generator additionally satisfies \(\Delta^2=0\).
Negative cohomology is zero, so \(\Delta(1)=0\). This is the convention of
the original question \cite[p.~2702]{OWR}; in degree one the bracket is the
commutator of derivations. Units are regarded as squarefree polynomials.

\begin{theorem}\label{main}
For every characteristic-zero field \(k\) and every nonzero \(h\in k[x]\),
the following are equivalent:
\begin{enumerate}
\item \(h\) is squarefree;
\item the canonical Gerstenhaber bracket on \(\HH^\bullet(A_h)\) is zero;
\item \(\HH^\bullet(A_h)\) admits a bracket generator;
\item \(\HH^\bullet(A_h)\) admits a BV generator.
\end{enumerate}
When \(h\) is squarefree, \(\Delta=0\) is a BV generator.
When \(h\) has a repeated irreducible factor, no bracket generator exists.
\end{theorem}

For \(h\ne0\), this presentation is the subalgebra of the first Weyl algebra
\(k\langle x,y\rangle/([y,x]-1)\) generated by \(x\) and \(Y=yh(x)\).
The original report uses \(h(x)y\); since \(yh-hy=h'\), the embedded
subalgebras agree. Lopes asks whether there are BV structures when this
family's Nakayama automorphism is not semisimple \cite[Question~(2)]{OWR}.
The theorem classifies all nonzero parameters, including every nonconstant
parameter in that branch. Section~\ref{scope} treats the literal embedded
specialization \(h=0\) separately.

The structural results used below are due to Benkart--Lopes--Ondrus
\cite{BLO} and Lopes--Solotar \cite{LS}. Nonexistence of a generator for
ordinary Hochschild Gerstenhaber cohomology is already known:
Herscovich--Li prove it for \(\operatorname{FK}(3)\)
\cite[Proposition~4.15]{HL}. Their use of cup/bracket relations to contradict
every candidate operator is an essential predecessor of the obstruction
strategy here. Our contribution is the squarefree classification for this
polynomial family, deduced from its established structure and the cup
comparison below. We make no claim to the first no-generator example or
an exhaustive historical priority determination.

\section{Published structure and elementary consequences}
Set \(A=A_h\), \(d=\gcd(h,h')\), with monic gcd, and \(\pi=h/d\).
The PBW basis is \(\{x^iY^j:i,j\ge0\}\). We use the following precise
published imports:
\begin{itemize}
\item The length-two free \(A\)-bimodule resolution of \cite[\S3,
(3.1)]{LS}, with degree-one generators \(x,Y\) and degree-two generator
\(r\), has
\begin{align}
\partial_2(r)
 ={}&[Y]x+Y[x]-[x]Y-x[Y]\notag\\[-2pt]
 &-\sum_{s\ge1}h_s\sum_{i=0}^{s-1}x^i[x]x^{s-1-i},
 \qquad h=\sum_s h_sx^s. \label{small}
\end{align}
Here \([x],[Y]\) denote free module generators, not commutator symbols.
\item In characteristic zero, \(\HH^2(A)=A/dA\) and
\(\HH^n(A)=0\) for \(n\ge3\)
\cite[Proposition~3.9 and Corollary~3.11]{LS}.
\item For \(f\in k[x]\), write \(D_f(x)=0,D_f(Y)=f\).
The classes of these derivations and of
\(\ad_{g a_n}\), \(g\in k[x]\), \(n\ge0\), span \(\HH^1(A)\).
In the localized ambient Weyl algebra,
\[
 a_n=\pi h^{n-1}y^n\quad(n\ge1),\qquad a_0=1/d.
\]
These elements normalize \(A\); moreover \(\ad_{g a_n}\) is inner
whenever \(d\mid g\). These are the derivation-classification results
of \cite[Theorems~5.3, 5.7, 5.13 and Proposition~5.9]{BLO},
explicitly recalled in \cite[\S6.1 and Theorem~6.2]{LS}.
We also use the ambient basis
\(A=\bigoplus_{m\ge0}k[x]h^m y^m\)
\cite[Proposition~2.1(b)]{LS}.
\end{itemize}
These are established structural inputs, not results certified by the
finite checks supplied with this note.

For completeness, the center and the needed innerness test are elementary.
If \(t=\sum_{j=0}^N t_j(x)Y^j\), \(N>0\), then the leading coefficient
of \([t,x]\) is \(Nh\,t_N\), of \(Y\)-degree \(N-1\).
Thus \(C_A(x)=k[x]\). Commuting also with \(Y\) gives \(ht'=0\), so
\(Z(A)=k=\HH^0(A)\). If \(D_f=\ad_t\), then \(t\in k[x]\) and
\(D_f(Y)=-ht'\). Conversely every multiple of \(h\) has that form by
polynomial integration. Consequently
\begin{equation}\label{inner}
 [D_f]=0\text{ in }\HH^1(A)\quad\Longleftrightarrow\quad h\mid f.
\end{equation}
We use \(\ad_t(z)=[t,z]\).

One can also see the stated quotient in the resolution directly.
The degree-one cochain differential is
\[
 (\alpha,\beta)\longmapsto
 [\beta,x]+[Y,\alpha]-\sum_s h_s\sum_{i=0}^{s-1}x^i\alpha x^{s-1-i}.
\]
The triangular PBW calculation gives \([A,x]=hA\), and coefficient
integration gives \([Y,A]=hA\). Modulo \(hA\), the last sum is
\(h'\alpha\). Therefore the image is \(hA+h'A=dA\).
The ideal \(dA\) is two-sided, since \([Y,d]=hd'\) is divisible by \(d\).

\begin{lemma}\label{central}
The class \(c=[D_d]\) is central in the Lie algebra \(\HH^1(A)\).
If \(d\ne1\), it is nonzero and has a linearly independent companion
\([D_{d/u}]\), for every irreducible \(u\mid d\).
\end{lemma}
\begin{proof}
The centrality also follows from \cite[Theorem~6.2(b)]{LS}; we give its
derivation from the spanning input to specify exactly what is used.
Extend \(D_d\) to the localized ambient Weyl algebra by
\(D_d(y)=1/\pi\), \(D_d(x)=0\). This respects \([y,x]=1\) and \(Y=yh\).
For \(f=1/\pi\), commuting \(y\) through coefficients gives
\[
 D_d(y^n)=\sum_{j=0}^{n-1}\binom{n}{j+1}f^{(j)}y^{n-1-j}.
\]
Define \(P_0=1\) and
\(P_{j+1}=\pi P_j'-(j+1)\pi'P_j\); induction gives
\(f^{(j)}=P_j/\pi^{j+1}\). Hence
\[
 D_d(g a_n)=
 \sum_{j=0}^{n-1}\binom{n}{j+1}
       g d^jP_j\,h^{n-1-j}y^{n-1-j}\in A\quad(n\ge1).
\]
Thus \([D_d,\ad_{g a_n}]=\ad_{D_d(g a_n)}\) is inner.
For \(n=0\), \(D_d(g/d)=0\); also \([D_d,D_f]=0\).
The spanning input proves centrality for all of \(\HH^1\).
If \(u\mid d\), then \(d\) and \(d/u\) are linearly independent
polynomials of degree less than \(\deg h\). Equation~\eqref{inner}
proves the last assertions.
\end{proof}

\section{The required canonical cup product}
\begin{lemma}\label{cup}
For every derivation \(E\) of \(A\) and every \(f\in k[x]\),
\[
 [D_f]\smile[E]=[\,fE(x)\,]\ \in\ \HH^2(A)=A/dA.
\]
In particular \(c\smile z=0\) for every \(z\in\HH^1(A)\).
\end{lemma}
\begin{proof}
Use normalized bars, with
\(b[a|b]=a[b]-[ab]+[a]b\).
A comparison from the small resolution to the normalized bar resolution
sends the degree-one generators to \([x],[Y]\), and \(r\) to
\begin{equation}\label{bar}
 F_2(r)=[Y|x]-[x|Y]
       -\sum_{s\ge2}h_s\sum_{i=1}^{s-1}[x^i|x]x^{s-1-i}.
\end{equation}
Indeed, for \(s\ge2\), telescoping gives
\[
 b\!\left(\sum_{i=1}^{s-1}[x^i|x]x^{s-1-i}\right)
 =\sum_{i=0}^{s-1}x^i[x]x^{s-1-i}-[x^s].
\]
The first two bars of \eqref{bar} contribute the four outer \(x,Y\)
terms and \(-[Yx]+[xY]=-[h]\). Constants vanish in normalized bars;
the linear term left over is \(-h_1[x]\).
This is precisely \eqref{small}. The degree-zero and degree-one chain
conditions are immediate, so this is a comparison lifting the identity.
Since the small complex is the published free resolution, pullback of
the ordinary bar cup product computes the canonical product.

For arbitrary derivations \(D,E\), its value on \(r\) is
\[
 D(Y)E(x)-D(x)E(Y)
 -\sum_{s\ge2}h_s\sum_{i=1}^{s-1}
       D(x^i)E(x)x^{s-1-i}.
\]
The right factor must remain on the right: \(E(x)\) need not commute
with \(x\). With \(D=D_f\), all polynomial correction terms and the
second term vanish literally, leaving \(fE(x)\).
For \(f=d\), this belongs to \(dA\), proving the last assertion.
\end{proof}

\section{An obstruction for all possible generators}
\begin{lemma}\label{linear}
Let \(H\) be a nonnegatively graded Gerstenhaber algebra over \(k\)
with \(H^0=k\). Suppose \(0\ne c\in H^1\) is Lie-central,
\(cH^1=0\), and \(\dim_k H^1>1\).
Every degree-minus-one bracket generator on \(H\) vanishes on \(H^1\).
Consequently every cup-zero pair in \(H^1\) would have zero bracket.
\end{lemma}
\begin{proof}
Write \(\lambda=\Delta|_{H^1}:H^1\to k\).
The degree-one part of \eqref{bv} is
\[
 [a,b]=-\Delta(ab)+\lambda(a)b-\lambda(b)a.
\]
Taking \(a=c\) gives \(\lambda(c)b=\lambda(b)c\).
A \(b\) independent of \(c\) first forces \(\lambda(c)=0\).
Then \(c\ne0\) forces \(\lambda(b)=0\) for every \(b\).
The displayed identity proves the consequence. No condition on
\(\Delta|_{H^2}\) or on \(\Delta^2\) was used.
\end{proof}

\begin{proof}[Proof of Theorem~\ref{main}]
If \(d=1\), the published spanning statement makes every
\(\ad_{g a_n}\) inner. Thus \(\HH^1(A)\) is spanned by the commuting
\([D_f]\). All brackets with \(\HH^0(A)=k\) vanish, and
\(\HH^{\ge2}(A)=0\). Therefore the entire bracket is zero and
\(\Delta=0\) is a BV generator.

Suppose \(d\ne1\). Choose an irreducible nonconstant \(u\mid d\)
over \(k\) itself, and set \(g=d/u\), \(a=[D_g]\).
If \(h=\alpha\prod_i u_i^{e_i}\) is its irreducible factorization, then
\(d=\prod_i u_i^{e_i-1}\) in characteristic zero, so
\(\pi=\alpha\prod_i u_i\).
In particular \(u\mid\pi\) and
\[
 L=\pi\,\frac{d'}d
   =\sum_i(e_i-1)\frac{\pi}{u_i}u_i'\ \in k[x].
\]
Define \(b\) directly on the presentation by
\(b(x)=\pi,\ b(Y)=LY\).
Write \(\beta=[b]\). It is a derivation because its value on the defining relation is
\[
 [LY,x]+[Y,\pi]-\pi h'
 =Lh+h\pi'-h'\pi=0.
\]
Lemma~\ref{cup} gives \(a\smile\beta=[g\pi]=0\), since
\(g\pi=d(\pi/u)\).
On generators the commutator has \(x\)-component zero and
\[
 [D_g,b](Y)=Lg-\pi g'
   =\frac{h}{u^2}u'=:q.
\]
Since \(u^2\mid h\) and \(u'\ne0\), this is a nonzero polynomial.
In characteristic zero,
\(\deg q=\deg h-\deg u-1<\deg h\), so \eqref{inner} gives
\([a,\beta]=[D_q]\ne0\).
Lemmas~\ref{central}--\ref{linear} now contradict every bracket
generator. They also exhibit a nonzero canonical bracket.
Finally (4) implies (3), and a zero bracket has the generator zero,
which supplies all stated equivalences.
\end{proof}

For \(h=x^n,\ n\ge2\), write \(\beta=[b]\). The pair is
\(a=[D_{x^{n-2}}]\), \(b(x)=x,\ b(Y)=(n-1)Y\), with
\([a,\beta]=a\ne0\) and \(a\smile\beta=0\).
For \(h=(x^2+1)^2\) over \(\mathbb Q\), take
\(a=[D_1]\), \(b(x)=x^2+1,\ b(Y)=2xY\);
the bracket of the two classes is \([D_{2x}]\ne0\).
No splitting field or scalar-extension claim is needed.
One should not identify \(b\) literally with \(\ad_{\pi y}\) in the
ambient convention \(Y=yh\): that latter derivation has the extra
polynomial term \(\pi h''-Lh'\) on \(Y\).
Our direct definition avoids that unnecessary identification.

\section{Original scope, predecessors, and verification}\label{scope}
The recorded Nakayama automorphism is
\(\nu(x)=x,\ \nu(Y)=Y+h'\)
\cite[p.~1381]{LS}. For nonconstant \(h\), its restriction to
\(\operatorname{span}_k\{Y,h'\}\) is a nontrivial unipotent Jordan block,
so it is not semisimple. For nonzero constant \(h\), it is the identity.
Thus Theorem~\ref{main} gives both positive and negative answers in the
whole original nonsemisimple branch; it is not an extension of a
semisimple-Nakayama theorem by omission of a hypothesis.

If \(h=0\) in the original \emph{embedded} definition, the generated
algebra is \(k[x]\). Here
\(\HH^0=k[x]\), \(\HH^1=k[x]\partial_x\), and higher cohomology vanishes.
Set \(\Delta(f\partial_x)=-f'\), \(\Delta(k[x])=0\).
The Leibniz rule verifies \eqref{bv} in degrees \((0,1)\) and \((1,1)\),
and the other cases vanish; \(\Delta^2=0\).
This separately answers that degenerate specialization.
Setting \(h=0\) in the \emph{abstract two-generator presentation} gives
a polynomial plane, which is a different algebra.

Kowalzig--Kr\"ahmer's ordinary-cohomology theorem
\cite[Theorem~1.7]{KK} assumes semisimplicity of the modular automorphism;
its construction uses twisted-coefficient homology in a duality.
It does not classify the present nonsemisimple branch.
Herscovich--Li's earlier no-generator theorem \cite{HL} concerns the
finite-dimensional three-generator Fomin--Kirillov algebra, not this
Ore-domain family, and already uses the stronger notion of generator
without square-zero. The recent positive theorem for truncated basic
cycle algebras \cite[Theorem~7.11]{BIKLZ} concerns another family;
its Criterion~7.7 includes commuting odd generators, unlike the
obstruction pair here. Neither theorem is contradicted or extended to
an unstated class.

Our bounded literature investigation through 22 October 2026 found no
authenticated prior statement of this complete \(A_h\) classification.
The structural classification and the general no-generator phenomenon
are inherited, and no exhaustive novelty or current-openness certificate
is asserted. We do not treat a dependency DOI as a prior publication of
this theorem. The separate general-algebra and Poisson questions in
\cite{OWR} are not claimed answered.

The source supplement includes a portable standard-library exact checker
for representative rational-polynomial identities and normalized-bar
telescoping, with deliberately false controls. It is newly written
supporting diagnostics, not execution of the original claimant's code
and not a proof of the universal theorem. The proof is the argument
above, subject to the explicitly named published structural inputs.
No formal proof-assistant verification is claimed.

\paragraph{Disclosure and license.}
AI tools were extensively involved in mathematical exploration, proof
development and reconstruction, computational verification, literature
and source investigation, manuscript preparation, editing, and
adversarial auditing. Alec Kriebel is the named human author.
This is an unrefereed preprint and has not undergone formal independent
human peer review. Internal AI reviews do not provide human refereeing
or journal acceptance. Original prose is licensed under CC BY 4.0;
original supporting code is MIT licensed.

\begin{thebibliography}{9}
\bibitem{BLO}
G.~Benkart, S.~A.~Lopes, and M.~Ondrus,
\emph{Derivations of a parametric family of subalgebras of the Weyl algebra},
J.~Algebra \textbf{424} (2015), 46--97.
\href{https://doi.org/10.1016/j.jalgebra.2014.11.007}{doi:10.1016/j.jalgebra.2014.11.007}.
\bibitem{BIKLZ}
X.~Bian, T.~Itagaki, W.~Kou, W.~Lyu, and G.~Zhou,
\emph{The Hochschild cohomlogy ring of a self-injective Nakayama
algebra is a Batalin--Vilkovisky algebra},
preprint, arXiv:2603.04834v2 (2026).
\href{https://arxiv.org/abs/2603.04834v2}{arxiv.org/abs/2603.04834v2}.
\bibitem{HL}
E.~Herscovich and Z.~Li,
\emph{Gerstenhaber structure on Hochschild cohomology of the
Fomin--Kirillov algebra on 3 generators},
Doc.~Math. \textbf{27} (2022), 1773--1804.
\href{https://doi.org/10.4171/DM/X18}{doi:10.4171/DM/X18}.
\bibitem{KK}
N.~Kowalzig and U.~Kr\"ahmer,
\emph{Batalin--Vilkovisky structures on Ext and Tor},
J.~Reine Angew.~Math. \textbf{697} (2014), 159--219.
\href{https://doi.org/10.1515/crelle-2012-0086}{doi:10.1515/crelle-2012-0086}.
\bibitem{OWR}
S.~A.~Lopes (joint work with A.~Solotar),
\emph{Gerstenhaber algebra structure on Hochschild cohomology},
in Mini-Workshop: Poisson and Poisson-type algebras,
Oberwolfach Report \textbf{46/2023}, 2702--2703, Question~(2).
Workshop 15--20 October 2023; report published 23 July 2024.
\href{https://doi.org/10.4171/OWR/2023/46}{doi:10.4171/OWR/2023/46}.
\bibitem{LS}
S.~A.~Lopes and A.~Solotar,
\emph{Lie structure on the Hochschild cohomology of a family of
subalgebras of the Weyl algebra},
J.~Noncommut.~Geom. \textbf{15} (2021), 1373--1407.
\href{https://doi.org/10.4171/JNCG/439}{doi:10.4171/JNCG/439}.
\end{thebibliography}
\end{document}
